\documentclass{article}%
\usepackage[T1]{fontenc}%
\usepackage[utf8]{inputenc}%
\usepackage{lmodern}%
\usepackage{textcomp}%
\usepackage{lastpage}%
\usepackage{geometry}%
\geometry{margin=2.2cm}%
%
\usepackage{amsmath}%
\usepackage{amssymb}%
\usepackage{booktabs}%
\usepackage{longtable}%
\usepackage{array}%
\usepackage{hyperref}%
\hypersetup{hidelinks}%
\title{ITA 2027 {-} Física (claude\_sonnet\_5\_5) x Poliedro\\[4pt]\large Comparação com o gabarito Poliedro Resolve (ITA 2027, 2ª fase, dia 3 {-} Física)\thanks{Gerado com gpt-resolve: \url{https://github.com/lgabs/gpt-resolve}}}%
\date{09/11/2026}%
%
\begin{document}%
\normalsize%
\maketitle%
\section*{Resumo}\textbf{Nota média: 9,5/10} (23 de 24 itens com a resolução correta; cada item vale o mesmo dentro da questão).\par\medskip%
\begin{longtable}{@{}c c p{0.78\textwidth}@{}}\toprule Questão & Nota & Resumo \\ \midrule
1 & 10,0 & A resolução está totalmente de acordo com o gabarito do Poliedro nos três itens, e os dois estão corretos. No item a, k\_c = MgL/2. No item b, a classificação é a mesma: r $\geq$ 1 dá só $\theta$ = 0, estável; 2/$\pi$ < r < 1 dá $\theta$ = 0 instável e $\pm$$\theta$$_{0}$ estáveis; r $\leq$ 2/$\pi$ dá só $\theta$ = 0, instável. No item c, a expressão de U\_total é a mesma e os três esboços têm a mesma forma: poço simples, poço duplo e máximo único. Refiz as contas com python3 (raiz $\theta$$_{0}$ $\approx$ 1,131 rad para r = 0,8, sinal de U'' em cada equilíbrio e valores de U nas bordas) e não encontrei erro de conta nem de física. Ficam só observações menores: o caso marginal r = 1 e o gráfico do gabarito, que não pude ver. \\
2 & 10,0 & Questão 2 correta nos dois itens e idêntica ao gabarito do Poliedro: a) H\_n = 0,49\^{}n V$^2$/(8g); b) $\Delta$t = V/{[}g(1-e){]} = 10V/(3g). Conferi o enunciado no PDF original. Refiz as contas (conservação de v\_x, v\_y multiplicada por e a cada choque, soma em PG dos tempos de voo) e conferi numericamente com python. O único ponto a registrar é que a palavra \textbackslash{}"energia\textbackslash{}" no item b é ambígua: lida como energia cinética, o mínimo seria atingido já no topo do primeiro voo, em t = V/(2g). Resolução e gabarito adotam a mesma leitura (energia mecânica), que é a pretendida pela questão. \\
3 & 10,0 & A resolução está certa nos dois itens, e as respostas finais são algebricamente iguais às do gabarito do Poliedro. Elas só mudam de forma: a resolução escreve as raízes como $\surd$($\mu$e$^2$g$^2$ - $\alpha$$^2$R$^2$), e o gabarito, como R$^2$$\cdot$$\surd$($\mu$e$^2$g$^2$/R$^2$ - $\alpha$$^2$). Conferi numericamente: com g = 9,8, $\mu$e = 0,5, $\mu$c = 0,3, R = 2, $\alpha$ = 0,5 e h = 1,2, as duas dão D = 2,5198 m e t\_parada = 4,6458 s. Esses valores batem com o cálculo pelas coordenadas explícitas do ponto de queda. O único problema está no gabarito: a resolução detalhada tem erros de digitação nas etapas intermediárias (falta o quadrado em $\mu$e$\cdot$g e faltam parênteses na última linha do item b), mas as respostas resumidas por item estão corretas. \\
4 & 10,0 & Questão correta nos dois itens. No item a, a resolução e o Poliedro dão a mesma resposta, h\_f = h\_i(1+M/m)\^{}\{3/5\}. No item b, a equação da resolução, 3(m+M)h\_i{[}1-(h\_i/h\_max)\^{}\{2/3\}{]} = 2m(h\_max-h\_i), é algebricamente a mesma da interpretação 1 do Poliedro (adiabática reversível). Conferi a equivalência numericamente e confirmei o valor de h\_max integrando a equação de movimento do pistão (2,2559 h\_i para m = M). O único ponto em aberto é a leitura do item b, que o próprio Poliedro marcou como ambígua. A escolha da resolução tem boa justificativa física e é a mais coerente com o pedido de uma \textbackslash{}"equação que relaciona\textbackslash{}" h\_max, m, M e h\_i. \\
5 & 5,0 & No item a, a resolução e o gabarito coincidem exatamente (p' e posições verticais $\pm$n1Rh/(n1R+(n2-n1)p)), e o traçado numérico de raios confirma. No item b, as duas fazem a mesma física com contas corretas, mas a resposta principal difere por um fator n1. A resolução toma $\lambda$ como o comprimento de onda no meio n1 e obtém n1$\lambda$D/(2n2h); o gabarito toma $\lambda$ no vácuo e obtém $\lambda$D/(2n2h). A divergência vem de uma ambiguidade do enunciado, não de erro de cálculo. A leitura do gabarito segue a convenção mais comum, e a resolução apresenta essa forma como alternativa explícita. Fica para o professor decidir se a resposta principal da resolução seria aceita integralmente. \\
6 & 10,0 & Os três itens da resolução coincidem com o gabarito do Poliedro: a) q = 3C$\varepsilon$/4; b) compressão de 2d/3, com distância final d/3; c) k = 27C$\varepsilon$$^2$/(64d$^2$). Refiz as contas com frações exatas. Também conferi o balanço de energia (dissipação de 9/32$\cdot$C$\varepsilon$$^2$) e a fatoração da força líquida feitos na resolução. Está tudo correto. A resolução é mais completa que o gabarito: declara as hipóteses (efeito pendular desprezado) e mostra que V3 = $\varepsilon$/4 é exatamente o limite de estabilidade do equilíbrio. Esse detalhe fisicamente curioso do enunciado não altera as respostas. Questão inteiramente correta. \\
7 & 10,0 & A resolução está correta nos dois itens e coincide com o gabarito Poliedro. Em a), |V\_o| = (V\_s/4)|G$\varepsilon$+$\alpha$$\Delta$T|. Em b), a ponte completa com quatro extensômetros (R2 e R3 em tração, R1 e R4 em compressão, ou o inverso) dá V\_o = V\_sG$\varepsilon$ sem efeito térmico. As expressões foram conferidas numericamente com o circuito exato. Não há divergência. \\
8 & 10,0 & A resolução da questão 8 está correta e coincide com o gabarito do Poliedro nos três itens. a) v = $\surd$(2qV/m). b) Movimento helicoidal uniforme (MU ao longo de B e MCU no plano perpendicular), bem justificado. c) cos$\theta$ = (LB/2$\pi$n)$\surd$(q/2mV), com n inteiro positivo. Conferi o enunciado no PDF original e verifiquei numericamente em Python que as expressões são equivalentes. A única diferença é a desigualdade estrita na restrição de n (cos$\theta$ < 1 contra cos$\theta$ $\leq$ 1), que só afeta o caso $\theta$ = 0 e não altera a nota. \\
9 & 10,0 & A resolução está totalmente de acordo com o gabarito do Poliedro nos três itens. No item a, ambas dão E = E0{[}1+(-1)\^{}n{]}cos(kx-$\omega$t)$\hat z$, uma onda que se propaga em +x. No item b, n par dá interferência construtiva (amplitude 2E0) e n ímpar dá destrutiva total. No item c, ambas dão B = -B0 cos(kx-$\omega$t)$\hat y$, com E$\times$B no sentido +x?. Refiz a dedução pela lei de Faraday e confere. Os acréscimos da resolução (B0 = 2E0/c, d = n$\lambda$/2, intensidade 4 vezes maior) estão corretos e não entram em conflito com o gabarito. Não há erro em nenhum dos dois, só observações de redação e de interpretação. \\
10 & 10,0 & A resolução está correta nos dois itens e confere com o gabarito do Poliedro. a) A fórmula de Compton é idêntica e a dedução também. b) As respostas são equivalentes: $\beta$ = ($\nu$$^2$ - $\nu$'$^2$)/($\nu$$^2$ + $\nu$'$^2$) = $\varepsilon$(2+$\varepsilon$)/(2+2$\varepsilon$+$\varepsilon$$^2$) = {[}(1+$\varepsilon$)$^2$ - 1{]}/{[}(1+$\varepsilon$)$^2$ + 1{]}, o que conferi algebricamente e numericamente. O sentido físico também coincide: O vai ao encontro do fóton espalhado. A única diferença é de convenção: a resolução escreve V com sinal negativo, o gabarito dá o módulo e diz o sentido por extenso. O gabarito tem erros de digitação só na dedução, não na resposta final. \\
\bottomrule\end{longtable}%
\section*{Principais pontos de dúvida}%
\begin{itemize}%
\item%
\textbf{Q4} \textit{(statement)} No item b, o próprio Poliedro diz que o enunciado não basta para achar p\_f em h\_max e propõe duas leituras. A resolução seguiu a 1 (adiabática reversível), que bate com o pedido de "equação que relaciona h\_max e m, M e h\_i". O gabarito oficial do ITA ainda não saiu. Se a banca esperar outra leitura, por exemplo a altura de equilíbrio final h\_eq = (1+3M/5m)h\_i, a nota do item b pode mudar. Vale levar a um professor.%
\item%
\textbf{Q5} \textit{(statement)} O enunciado diz apenas que as fontes têm "comprimento de onda $\lambda$" e estão no meio n1, sem dizer se $\lambda$ vale no vácuo ou em n1. Isso altera a resposta do item b por um fator n1: $\lambda$D/(2n2h) se for no vácuo, n1$\lambda$D/(2n2h) se for em n1. O gabarito Poliedro adota o vácuo e a resolução adota n1 como resposta principal, citando a outra como alternativa. Vale levar ao professor para saber qual leitura a banca deve aceitar e se a resposta principal da resolução levaria desconto.%
\end{itemize}%
\section*{Questão 1}%
\href{https://poliedroresolve.sistemapoliedro.com.br/vestibulares/ita/2027/ita-2-fase-dia-3-22-10-2026-2-fase-dia-3/dissertativa-1-fisica-2-fase-dia-3-ita-2027}{Resolução do gabarito}\par\smallskip%
\begin{longtable}{@{}p{0.06\textwidth} p{0.33\textwidth} p{0.33\textwidth} p{0.1\textwidth} p{0.08\textwidth}@{}}\toprule Item & Resolução & Gabarito & Veredito & Certo \\ \midrule
a & k\_c = MgL/2 & k\_c = MgL/2 & igual & ambos \\
b & r = k/k\_c $\geq$ 1: só $\theta$ = 0, estável; 2/$\pi$ < r < 1: três equilíbrios, $\theta$ = 0 instável e $\pm$$\theta$$_{0}$ estáveis (sen $\theta$$_{0}$/$\theta$$_{0}$ = r); 0 < r $\leq$ 2/$\pi$: só $\theta$ = 0, instável & k/k\_c $\geq$ 1: um equilíbrio ($\theta$ = 0), estável; 2/$\pi$ < k/k\_c < 1: três equilíbrios, $\theta$ = 0 instável e $\pm$$\theta$$_{0}$ estáveis (k$\theta$$_{0}$ = k\_c sen $\theta$$_{0}$); k/k\_c $\leq$ 2/$\pi$: um equilíbrio ($\theta$ = 0), instável & igual & ambos \\
c & U = (MgL/2)cos $\theta$ + k$\theta$$^2$/2. Esboços: poço simples em 0 (r $\geq$ 1); poço duplo simétrico, com máximo local em 0 e mínimos em $\pm$$\theta$$_{0}$ (2/$\pi$ < r < 1); máximo único em 0, com U decrescente em |$\theta$| (r $\leq$ 2/$\pi$) & U($\theta$) = (MgL/2)cos $\theta$ + k$\theta$$^2$/2. k $\geq$ k\_c: poço único com mínimo em 0; 2/$\pi$ < k/k\_c < 1: duplo poço (máximo local em 0, mínimos em $\pm$$\theta$$_{0}$); k/k\_c $\leq$ 2/$\pi$: máximo em 0 e U decrescendo para as bordas & igual & ambos \\
\bottomrule\end{longtable}%
\begin{itemize}%
\item%
\textbf{a)} Com U = (MgL/2)cos$\theta$ + k$\theta$$^2$/2, U''(0) = k - MgL/2, que muda de sinal em k = MgL/2. Para $\theta$ ? 1, o torque resultante fica (MgL/2 - k)$\theta$. A resolução e o gabarito chegam a esse resultado pelo mesmo raciocínio. A unidade (N$\cdot$m/rad) está correta.%
\item%
\textbf{b)} As faixas de r, o número de equilíbrios e a natureza de cada um coincidem. A condição sen $\theta$$_{0}$/$\theta$$_{0}$ = r equivale a k$\theta$$_{0}$ = k\_c sen $\theta$$_{0}$. Conferi os detalhes: sen $\theta$/$\theta$ é estritamente decrescente em (0, $\pi$/2), de 1 a 2/$\pi$, e por isso a raiz é única. U''($\theta$$_{0}$) = (k\_c/$\theta$$_{0}$)(sen $\theta$$_{0}$ - $\theta$$_{0}$cos $\theta$$_{0}$) > 0, então $\pm$$\theta$$_{0}$ são estáveis. Em r = 1, U $\approx$ k\_c(1 + $\theta$$^{4}$/24), um mínimo (estável). Em r = 2/$\pi$ a raiz cairia em $\pm$$\pi$/2, que está fora do intervalo aberto. Exemplo numérico: com r = 0,8 sai $\theta$$_{0}$ $\approx$ 1,131 rad $\approx$ 64,8$^\circ$, que bate com a resolução.%
\item%
\textbf{c)} A expressão de U e a forma qualitativa dos três gráficos são as mesmas. A referência de altura no mancal (U(0) = MgL/2) foi respeitada. Os valores dos gráficos da resolução estão certos: com r = 1,5, U($\pi$/2)/k\_c $\approx$ 1,85; com r = 0,8, U($\theta$$_{0}$)/k\_c $\approx$ 0,937 e U($\pi$/2)/k\_c $\approx$ 0,987; com r = 0,4, U($\pi$/2)/k\_c $\approx$ 0,493. Todos cabem nos eixos escolhidos. Não consegui ver a figura do gabarito para o item c (é uma imagem), mas a descrição em texto coincide.%
\end{itemize}%
\paragraph{Pontos de dúvida}%
\begin{itemize}%
\item%
\textit{(statement)} No caso{-}limite k = k\_c, U''(0) = 0 e a estabilidade sai do termo de 4ª ordem (U $\approx$ k\_c(1 + $\theta$$^{4}$/24)). O enunciado diz que $\theta$ = 0 deixa de ser estável 'para k < k\_c', o que implica estável em k = k\_c. A resolução e o gabarito seguem essa leitura, mas é um equilíbrio marginal (mínimo muito achatado).%
\item%
\textit{(key)} O gráfico do item c no gabarito é uma imagem que não pude ver. A conferência foi feita só pela descrição em texto, que coincide com a resolução.%
\item%
\textit{(solution)} A aproximação $\theta$$_{0}$ $\approx$ $\surd$(6(1-r)) só vale para r próximo de 1 (com r = 0,8 dá 1,095 contra o valor exato 1,131). A resolução avisa disso, e a aproximação não entra na resposta.%
\end{itemize}%
\section*{Questão 2}%
\href{https://poliedroresolve.sistemapoliedro.com.br/vestibulares/ita/2027/ita-2-fase-dia-3-22-10-2026-2-fase-dia-3/dissertativa-2-fisica-2-fase-dia-3-ita-2027}{Resolução do gabarito}\par\smallskip%
\begin{longtable}{@{}p{0.06\textwidth} p{0.33\textwidth} p{0.33\textwidth} p{0.1\textwidth} p{0.08\textwidth}@{}}\toprule Item & Resolução & Gabarito & Veredito & Certo \\ \midrule
a & H\_n = (0,7)\^{}\{2n\} V$^2$/(8g) = (0,49)\^{}n V$^2$/(8g) & h\_n = (e\^{}n V sen30$^\circ$)$^2$/(2g) = 0,7\^{}\{2n\} V$^2$/(8g) & igual & ambos \\
b & $\Delta$t = V/{[}g(1-e){]} = 10V/(3g) $\approx$ 3,3 V/g; nesse instante E\_min = (3/8)mV$^2$ e a massa desliza com v = ($\surd$3/2)V & T = (2V sen30$^\circ$/g)$\cdot$1/(1-e) = 10V/(3g) & igual & ambos \\
\bottomrule\end{longtable}%
\begin{itemize}%
\item%
\textbf{a)} A resolução e o gabarito usam o mesmo raciocínio. Como o solo é liso, v\_x se conserva e cada choque multiplica v\_y por e, então v\_\{y,n\} = e\^{}n$\cdot$V/2 e H\_n = v\_\{y,n\}$^2$/(2g) = 0,49\^{}n V$^2$/(8g). Conferi numericamente para n = 0 a 3. A resolução também confere o caso n = 0, que dá V$^2$/(8g), o resultado do lançamento oblíquo sem choque.%
\item%
\textbf{b)} Os dois interpretam a energia como energia mecânica total. O mínimo dessa energia é atingido depois de infinitos choques, mas em tempo finito, quando não resta velocidade vertical. Os tempos de voo formam uma PG: $\tau$\_k = 2e\^{}k(V/2)/g = e\^{}k V/g. A soma vale V/(0,3g) = 10V/(3g), o que conferi numericamente somando 200 termos. A resolução acrescenta o valor E\_min = (3/8)mV$^2$ e verifica os limites e$\to$0 e e$\to$1, ambos coerentes.%
\end{itemize}%
\paragraph{Pontos de dúvida}%
\begin{itemize}%
\item%
\textit{(statement)} O enunciado fala em "mínimo da sua energia" sem dizer qual energia. A resolução e o gabarito tomam a energia mecânica total, que só chega a (3/8)mV$^2$ depois de infinitos choques, em T = 10V/(3g). Se "energia" fosse lida como energia cinética, o valor mínimo (3/8)mV$^2$ já seria atingido pela primeira vez no topo do primeiro voo, em t = V/(2g), e se repetiria no topo de cada voo seguinte. Como o enunciado dá o coeficiente de restituição, a leitura pretendida é claramente a da energia mecânica, mas vale registrar a ambiguidade.%
\item%
\textit{(solution)} A resolução diz que o mínimo "só é atingido no limite". Na verdade, em t = T exato a componente vertical é nula e y = 0, então E = (3/8)mV$^2$ é de fato atingida nesse instante, e depois a massa desliza. É só imprecisão de redação; a resposta não muda.%
\item%
\textit{(key)} Na transcrição da resolução do gabarito faltam parênteses: aparece "(e\^{}nVsen$\theta$\^{}2)/2g" e "0,7\^{}2nV\^{}20,5\^{}2)/2gh". Isso parece erro de formatação do texto extraído, não de física, porque a resposta final 0,7\^{}\{2n\}V$^2$/(8g) está correta.%
\end{itemize}%
\section*{Questão 3}%
\href{https://poliedroresolve.sistemapoliedro.com.br/vestibulares/ita/2027/ita-2-fase-dia-3-22-10-2026-2-fase-dia-3/dissertativa-3-fisica-2-fase-dia-3-ita-2027}{Resolução do gabarito}\par\smallskip%
\begin{longtable}{@{}p{0.06\textwidth} p{0.33\textwidth} p{0.33\textwidth} p{0.1\textwidth} p{0.08\textwidth}@{}}\toprule Item & Resolução & Gabarito & Veredito & Certo \\ \midrule
a & d = $\surd$(R$^2$ + (2hR/g)$\cdot$$\surd$($\mu$e$^2$g$^2$ - $\alpha$$^2$R$^2$)) & D = R$\cdot$$\surd$(1 + (2h/g)$\cdot$$\surd$($\mu$e$^2$g$^2$/R$^2$ - $\alpha$$^2$)), com $\omega$\_e$^2$ = $\surd$($\mu$e$^2$g$^2$/R$^2$ - $\alpha$$^2$) & equivalente & ambos \\
b & t\_para = $\surd$(2h/g) + $\surd$R$\cdot$($\mu$e$^2$g$^2$ - $\alpha$$^2$R$^2$)\^{}\{1/4\}$\cdot${[}1/($\alpha$R) + 1/($\mu$c$\cdot$g){]} & t\_parada = $\omega$\_e/$\alpha$ + $\surd$(2h/g) + $\omega$\_e$\cdot$R/($\mu$c$\cdot$g), com $\omega$\_e = ($\mu$e$^2$g$^2$/R$^2$ - $\alpha$$^2$)\^{}\{1/4\} & equivalente & ambos \\
\bottomrule\end{longtable}%
\begin{itemize}%
\item%
\textbf{a)} Pondo R$^2$ em evidência: (2hR/g)$\cdot$$\surd$($\mu$e$^2$g$^2$ - $\alpha$$^2$R$^2$) = R$^2$$\cdot$(2h/g)$\cdot$$\surd$($\mu$e$^2$g$^2$/R$^2$ - $\alpha$$^2$). As duas formas são iguais. As duas resoluções seguem o mesmo caminho: o escorregamento começa quando $\mu$e$\cdot$g = R$\surd$($\alpha$$^2$ + $\omega$$^{4}$), daí vem um lançamento horizontal com v = $\omega$\_e$\cdot$R tangente à borda, e D$^2$ = R$^2$ + (v$\cdot$t\_q)$^2$, porque o raio é perpendicular à tangente. Conferi numericamente (g = 9,8; $\mu$e = 0,5; R = 2; $\alpha$ = 0,5; h = 1,2): as duas fórmulas dão 2,5198 m, o mesmo valor de $\surd$(x\_f$^2$ + y\_f$^2$) calculado com coordenadas explícitas.%
\item%
\textbf{b)} $\omega$\_e = ($\mu$e$^2$g$^2$/R$^2$ - $\alpha$$^2$)\^{}\{1/4\} = ($\mu$e$^2$g$^2$ - $\alpha$$^2$R$^2$)\^{}\{1/4\}/$\surd$R = Q/$\surd$R. Portanto $\omega$\_e/$\alpha$ = Q/($\alpha$$\surd$R) = $\surd$R$\cdot$Q/($\alpha$R) e $\omega$\_e$\cdot$R/($\mu$c$\cdot$g) = Q$\surd$R/($\mu$c$\cdot$g), que são exatamente os termos da resolução. As duas respostas somam o tempo até escorregar, o tempo de queda e o tempo de frenagem com desaceleração $\mu$c$\cdot$g, supondo que a componente horizontal da velocidade se conserva no contato com o solo. No exemplo numérico, as duas dão 4,6458 s.%
\end{itemize}%
\paragraph{Pontos de dúvida}%
\begin{itemize}%
\item%
\textit{(key)} A resolução detalhada do gabarito tem erros de digitação nas etapas intermediárias: aparece '$\omega$\_e$^2$ = $\surd$($\mu$\_e g/R - $\alpha$$^2$)' e 'D = R$\surd$(1 + 2h/g$\cdot$$\surd$($\mu$\_e g/R$^2$ - $\alpha$$^2$))', sem o quadrado em $\mu$\_e$\cdot$g (e, na primeira, também sem o R$^2$). A última linha, 't\_P = 1/$\alpha$ + R/$\mu$\_c g $^{4}$$\surd$(...)', também está sem parênteses: o certo é (1/$\alpha$ + R/($\mu$c$\cdot$g))$\cdot$$\omega$\_e + $\surd$(2h/g). As respostas resumidas por item estão corretas, mas quem for usar a resolução detalhada deve se guiar pelas fórmulas com $\mu$e$^2$g$^2$/R$^2$.%
\item%
\textit{(statement)} O enunciado não informa que $\mu$e$\cdot$g > $\alpha$$\cdot$R. Sem essa condição, a massa escorregaria já em t = 0 e as raízes não seriam reais. As duas resoluções supõem isso, mas só a resolução conferida diz isso explicitamente.%
\item%
\textit{(statement)} No item b, as duas resoluções supõem que, ao tocar o solo sem quicar, a velocidade vertical é anulada sem impulso de atrito, e a velocidade horizontal v = $\omega$\_e$\cdot$R se conserva. Se o atrito agisse de forma impulsiva no choque, o resultado mudaria; a resolução comenta isso numa ressalva. Essa leitura ('realiza um deslizamento') é a usada pelo gabarito.%
\item%
\textit{(solution)} As duas resoluções supõem que a massa sai do disco no instante em que começa a escorregar, com a velocidade da borda. A justificativa é que a aceleração relativa tem componente radial para fora e a massa já está na borda. A resolução justifica isso; o gabarito, não.%
\end{itemize}%
\section*{Questão 4}%
\href{https://poliedroresolve.sistemapoliedro.com.br/vestibulares/ita/2027/ita-2-fase-dia-3-22-10-2026-2-fase-dia-3/dissertativa-4-fisica-2-fase-dia-3-ita-2027}{Resolução do gabarito}\par\smallskip%
\begin{longtable}{@{}p{0.06\textwidth} p{0.33\textwidth} p{0.33\textwidth} p{0.1\textwidth} p{0.08\textwidth}@{}}\toprule Item & Resolução & Gabarito & Veredito & Certo \\ \midrule
a & h\_f = h\_i (1 + M/m)\^{}\{3/5\} (adiabática reversível, PV\^{}\{5/3\} = cte, com P\_i = (m+M)g/A e P\_f = mg/A) & h\_f = (1 + M/m)\^{}\{3/5\} $\cdot$ h\_i (adiabática reversível, $\gamma$ = 5/3) & igual & ambos \\
b & 3(m+M) h\_i {[}1 - (h\_i/h\_max)\^{}\{2/3\}{]} = 2m (h\_max - h\_i), com a raiz h\_max > h\_i. Equivale a 2m y\^{}4 + 2m y\^{}3 + 2m y\^{}2 - 3(m+M)(y+1) = 0, com y = (h\_max/h\_i)\^{}\{1/3\}. Hipóteses: energia conservada e gás em adiabática reversível. & Forma geral: h\_max = (3M+5m) g h\_i/(3 p\_f A + 2mg), com p\_f indeterminado pelo enunciado. Interpretação 1 (adiabática reversível): 3(M+m)(h\_i/h\_max)\^{}\{5/3\} h\_max + 2m h\_max = (3M+5m) h\_i. Interpretação 2 (repouso final após as oscilações): h\_eq = (1 + 3M/(5m)) h\_i. & equivalente & ambos \\
\bottomrule\end{longtable}%
\begin{itemize}%
\item%
\textbf{a)} Mesma expressão e mesmo raciocínio: a retirada quase estática com paredes e pistão isolantes é uma adiabática reversível. P\_i(Ah\_i)\^{}\{5/3\} = P\_f(Ah\_f)\^{}\{5/3\} e P\_i/P\_f = (m+M)/m dão h\_f = h\_i(1+M/m)\^{}\{3/5\}. Na transcrição do gabarito aparece "p\_i = (M + mg)/A", que é só erro de digitação de (M+m)g/A.%
\item%
\textbf{b)} A equação da resolução é algebricamente idêntica à da interpretação 1 do gabarito. Basta passar 3(m+M)h\_i(h\_i/h\_max)\^{}\{2/3\} para o outro lado, lembrando que (h\_i/h\_max)\^{}\{5/3\}$\cdot$h\_max = h\_i\^{}\{5/3\}h\_max\^{}\{{-}2/3\}. Conferi numericamente: sol + key $\approx$ 1e{-}14 para parâmetros aleatórios. O polinômio em y também está certo: sol$\cdot$y$^2$/h\_i = -(y-1)$\cdot$P(y). Integrei por RK4 a equação de movimento do pistão, m$\cdot$? = P\_i A (h\_i/h)\^{}\{5/3\} - mg. Com m = M o pistão para pela primeira vez em 2,2559 h\_i, igual à raiz da equação (2,2559). O h\_f $\approx$ 1,516 h\_i confere. Para M = 0,01m deu 1,0120 $\approx$ 1 + 6M/(5m). A física da resolução está certa. Sem atrito, a velocidade do pistão é contínua e muito menor que a do som no gás, então não há dissipação e a evolução é isentrópica (o pistão oscila sem amortecer). A conservação de energia $\Delta$U + mg$\Delta$h = 0 em v = 0 coincide com a 1ª lei usada pelo gabarito. A forma geral do gabarito também está correta, mas fica em função de p\_f. A interpretação 2 dá a altura de equilíbrio depois do amortecimento, não h\_max. Como o enunciado pede "a equação que relaciona h\_max e m, M e h\_i" (uma relação implícita), a leitura da resolução é a mais natural.%
\end{itemize}%
\paragraph{Pontos de dúvida}%
\begin{itemize}%
\item%
\textit{(statement)} No item b, o próprio Poliedro diz que o enunciado não basta para achar p\_f em h\_max e propõe duas leituras. A resolução seguiu a 1 (adiabática reversível), que bate com o pedido de "equação que relaciona h\_max e m, M e h\_i". O gabarito oficial do ITA ainda não saiu. Se a banca esperar outra leitura, por exemplo a altura de equilíbrio final h\_eq = (1+3M/5m)h\_i, a nota do item b pode mudar. Vale levar a um professor.%
\item%
\textit{(key)} O Poliedro diz que a relação de Poisson exige uma hipótese extra. No modelo ideal do enunciado (sem atrito, vácuo acima, pistão bem mais lento que o som no gás, velocidade inicial nula e sem choque), a evolução do gás é de fato isentrópica e o pistão oscila entre h\_i e h\_max sem amortecer. A irreversibilidade só aparece com dissipação. A ressalva do Poliedro é cautelosa, mas a hipótese da resolução tem boa justificativa física. Já a interpretação 2 do Poliedro não responde h\_max, e sim a altura final depois das oscilações amortecidas.%
\item%
\textit{(key)} A transcrição da resolução do gabarito tem erros tipográficos e de formatação que não mudam a resposta: "p\_i = (M + mg)/A" no lugar de (M+m)g/A, "h\_f = 1 + M/m\^{}(3/5) $\cdot$ h\_i" sem parênteses e h\_f usado no lugar de h\_max na 1ª lei do item b.%
\item%
\textit{(solution)} Detalhes de redação na resolução. A frase "Não há atrito nem vácuo a dissipar energia" ficou estranha. A aproximação h\_max $\approx$ 3(m+M)h\_i/(2m) para M ? m só vale assintoticamente e converge devagar: com M = 20m, o valor exato é 29,2 h\_i e a aproximação dá 31,5 h\_i. Nada disso afeta a resposta final.%
\end{itemize}%
\section*{Questão 5}%
\href{https://poliedroresolve.sistemapoliedro.com.br/vestibulares/ita/2027/ita-2-fase-dia-3-22-10-2026-2-fase-dia-3/dissertativa-5-fisica-2-fase-dia-3-ita-2027}{Resolução do gabarito}\par\smallskip%
\begin{longtable}{@{}p{0.06\textwidth} p{0.33\textwidth} p{0.33\textwidth} p{0.1\textwidth} p{0.08\textwidth}@{}}\toprule Item & Resolução & Gabarito & Veredito & Certo \\ \midrule
a & p' = -n2$\cdot$R$\cdot$p/(n1$\cdot$R + (n2-n1)$\cdot$p): imagem virtual no meio n1, a |p'| = n2Rp/(n1R+(n2-n1)p) do vértice. Posições verticais: y' = $\pm$n1$\cdot$R$\cdot$h/(n1$\cdot$R + (n2-n1)$\cdot$p), imagens direitas e menores. & p' = -n2$\cdot$p$\cdot$R/(n1$\cdot$R + (n2-n1)$\cdot$p): imagens virtuais no meio n1. Posições verticais: y' = $\pm$n1$\cdot$R$\cdot$h/(n1$\cdot$R + (n2-n1)$\cdot$p). & igual & ambos \\
b & y1 = n1$\cdot$$\lambda$$\cdot$D/(2$\cdot$n2$\cdot$h), com $\lambda$ tomado como o comprimento de onda no meio n1. A resolução diz explicitamente que, se $\lambda$ for o do vácuo, y1 = $\lambda$D/(2n2h). & y = $\lambda$$\cdot$D/(2$\cdot$n2$\cdot$h), com $\lambda$ tomado como o comprimento de onda no vácuo. & diferente & indefinido \\
\bottomrule\end{longtable}%
\begin{itemize}%
\item%
\textbf{a)} As expressões são idênticas. Refiz pela equação do dioptro com o centro de curvatura do lado da luz incidente (R\_dioptro = -R): n1/p + n2/p' = -(n2-n1)/R. O aumento é m = -n1p'/(n2p) = n1R/(n1R+(n2-n1)p). Conferi com um traçado exato de raios (n1=1; n2=1,5; R=10; p=5; h=0,01): imagem em x = -6,0000 e altura 0,0080, como dão as fórmulas (p' = -6; h' = 0,8h). A resolução deduz a fórmula diretamente da lei de Snell e o gabarito usa a equação do dioptro. O resultado é o mesmo.%
\item%
\textbf{b)} A física e as contas das duas versões estão corretas. A diferença, um fator n1, vem só de como cada uma lê o $\lambda$ do enunciado. Com R$\to$$\infty$, as imagens ficam em $\pm$h, a n2p/n1 da interface. Pelo princípio de Fermat, a diferença de caminho óptico entre as duas fontes até a altura y da tela é $\Delta$ = 2hy$\cdot$n1n2/(n1D+n2p) $\approx$ 2n2hy/D, já que p?D. A primeira franja clara exige $\Delta$ = $\lambda$\_vácuo, o que dá y = $\lambda$\_vácuo$\cdot$D/(2n2h). Conferi numericamente (n1=1,33; n2=1,6; p=1 cm; D=2 m; h=0,5 mm; $\lambda$=500 nm): o resultado exato foi 0,6288 mm e a fórmula aproximada dá 0,6250 mm. A diferença é só o termo de ordem p/D. Se $\lambda$ é o comprimento de onda no meio n1, então $\lambda$\_vácuo = n1$\lambda$ e chega{-}se exatamente à resposta principal da resolução. O gabarito adota o vácuo, que é a convenção mais usual para descrever uma fonte. A leitura da resolução também se sustenta, porque o enunciado dá $\lambda$ para fontes imersas em n1 sem dizer o meio. Além disso, a resolução declara a hipótese e apresenta a forma do gabarito como alternativa. Só coincidem se n1 = 1.%
\end{itemize}%
\paragraph{Pontos de dúvida}%
\begin{itemize}%
\item%
\textit{(statement)} O enunciado diz apenas que as fontes têm "comprimento de onda $\lambda$" e estão no meio n1, sem dizer se $\lambda$ vale no vácuo ou em n1. Isso altera a resposta do item b por um fator n1: $\lambda$D/(2n2h) se for no vácuo, n1$\lambda$D/(2n2h) se for em n1. O gabarito Poliedro adota o vácuo e a resolução adota n1 como resposta principal, citando a outra como alternativa. Vale levar ao professor para saber qual leitura a banca deve aceitar e se a resposta principal da resolução levaria desconto.%
\item%
\textit{(solution)} A resposta principal do item b pressupõe que $\lambda$ é o comprimento de onda no meio n1. Pela convenção mais usual ($\lambda$ da fonte = $\lambda$ no vácuo), a resposta seria $\lambda$D/(2n2h). A resolução cita essa forma só numa observação final. Um corretor rigoroso pode olhar apenas a resposta destacada.%
\item%
\textit{(statement)} O item a pede as posições verticais "em função de n1, n2, R e p", mas a resposta precisa também de h. Tanto a resolução quanto o gabarito incluem h, o que é o correto.%
\item%
\textit{(key)} No item b, o gabarito usa sen $\theta$2 $\approx$ y/D e despreza o recuo das imagens virtuais (a distância efetiva até a tela é D + n2p/n1). Isso é justificado por p ? D. O resultado exato em ordem paraxial é y = $\lambda$(n1D+n2p)/(2hn1n2), que tende à forma do gabarito quando p/D $\to$ 0.%
\end{itemize}%
\section*{Questão 6}%
\href{https://poliedroresolve.sistemapoliedro.com.br/vestibulares/ita/2027/ita-2-fase-dia-3-22-10-2026-2-fase-dia-3/dissertativa-6-fisica-2-fase-dia-3-ita-2027}{Resolução do gabarito}\par\smallskip%
\begin{longtable}{@{}p{0.06\textwidth} p{0.33\textwidth} p{0.33\textwidth} p{0.1\textwidth} p{0.08\textwidth}@{}}\toprule Item & Resolução & Gabarito & Veredito & Certo \\ \midrule
a & q = 3C$\varepsilon$/4 (V3 = $\varepsilon$/4, C' = 3C) & q = 3C$\varepsilon$/4 (C' = 3C, V3 = $\varepsilon$/4) & igual & ambos \\
b & compressão x = 2d/3 (distância final d/3) & compressão $\Delta$x = 2d/3 (distância final d/3) & igual & ambos \\
c & k = 27C$\varepsilon$$^2$/(64d$^2$) = 27$\varepsilon$0a$^2$$\varepsilon$$^2$/(64d$^3$) & k = 27C$\varepsilon$$^2$/(64d$^2$) & igual & ambos \\
\bottomrule\end{longtable}%
\begin{itemize}%
\item%
\textbf{a)} Refiz a conta. Em t2 o capacitor está com Q0 = C$\varepsilon$ e tensão $\varepsilon$, então V3 = $\varepsilon$/4. Conservando a carga, C$\varepsilon$ = (C + C')$\cdot$$\varepsilon$/4, o que dá C' = 3C, e a carga nas placas é q = C'V3 = 3C$\varepsilon$/4. A resolução chega a isso pela conservação de carga em cada conjunto isolado (q = C($\varepsilon$ - V3)) e o gabarito pela associação em paralelo. Os dois caminhos são equivalentes.%
\item%
\textbf{b)} C' = $\varepsilon$0A/d' = 3$\varepsilon$0A/d, logo d' = d/3. A mola, que estava no comprimento natural d, sofre compressão de d - d/3 = 2d/3. Os dois indicam corretamente que é compressão.%
\item%
\textbf{c)} F\_e = q$^2$/(2$\varepsilon$0A). Com $\varepsilon$0A = Cd, F\_e = (9C$^2$$\varepsilon$$^2$/16)/(2Cd) = 9C$\varepsilon$$^2$/(32d). Pelo equilíbrio, k$\cdot$(2d/3) = F\_e, logo k = 27C$\varepsilon$$^2$/(64d$^2$). Conferi com frações exatas. A resolução também escreve a forma em $\varepsilon$0a$^2$, que é equivalente porque C = $\varepsilon$0a$^2$/d. Conferi ainda, com frações exatas, o balanço de energia da resolução (energia final 7/32$\cdot$C$\varepsilon$$^2$ e dissipação 9/32$\cdot$C$\varepsilon$$^2$) e a fatoração da força líquida, -C$\varepsilon$$^2$(3s-d)$^2$(3s+5d)/{[}64d$^2$(s+d)$^2${]}.%
\end{itemize}%
\paragraph{Pontos de dúvida}%
\begin{itemize}%
\item%
\textit{(statement)} As placas estão penduradas por fios, e as instruções gerais da prova falam em sen $\theta$ $\approx$ $\theta$. Mesmo assim, o enunciado não dá massa das placas nem comprimento dos fios. Por isso a resolução e o gabarito desprezam a componente restauradora do peso (efeito pendular) e consideram só a força elétrica e a força da mola na horizontal. A hipótese é razoável e é a única que torna o problema solúvel. A resolução a declara; o gabarito a deixa implícita.%
\item%
\textit{(statement)} V3 = $\varepsilon$/4 cai exatamente no limite de estabilidade. Com o k obtido, a força líquida F(s) é $\leq$ 0 para qualquer s e tem raiz dupla em s = d/3. Ou seja, o equilíbrio é semiestável e só seria atingido assintoticamente, com amortecimento forte; para V3 < $\varepsilon$/4 não existe equilíbrio estável (s > d/3 é necessário). As respostas pedidas não mudam, mas fisicamente o cenário do enunciado é uma idealização. A resolução comenta isso; o gabarito não.%
\item%
\textit{(statement)} 'Tensão em t2' foi lida como $\varepsilon$, a tensão do capacitor carregado no instante da troca das chaves. Se alguém usasse $\varepsilon$/2 (tensão logo após fechar S2 com a distância ainda d), teria V3 = $\varepsilon$/8 e todos os resultados mudariam. A resolução e o gabarito adotam $\varepsilon$, que é a leitura natural. A resolução ainda descarta $\varepsilon$/8 porque levaria a um equilíbrio instável (s = d/7).%
\item%
\textit{(key)} A transcrição da resolução do gabarito tem fórmulas truncadas pela extração, por exemplo 'C$\varepsilon$ = C + C'$\varepsilon$/4C' = 3C' e 'k2d/3 = ...'. O sentido se reconstrói sem problema (C$\varepsilon$ = (C + C')$\varepsilon$/4 $\Rightarrow$ C' = 3C; k$\cdot$2d/3 = ...), e não há erro de física.%
\end{itemize}%
\section*{Questão 7}%
\href{https://poliedroresolve.sistemapoliedro.com.br/vestibulares/ita/2027/ita-2-fase-dia-3-22-10-2026-2-fase-dia-3/dissertativa-7-fisica-2-fase-dia-3-ita-2027}{Resolução do gabarito}\par\smallskip%
\begin{longtable}{@{}p{0.06\textwidth} p{0.33\textwidth} p{0.33\textwidth} p{0.1\textwidth} p{0.08\textwidth}@{}}\toprule Item & Resolução & Gabarito & Veredito & Certo \\ \midrule
a & |V\_o| = (V\_s/4)|G$\varepsilon$ + $\alpha$$\Delta$T|, para o extensômetro em qualquer braço (o sinal depende do braço: - em R1 e R4, + em R2 e R3) & |V0| = (Vs/4)$\cdot$|G$\varepsilon$ + $\alpha$$\Delta$T| & igual & ambos \\
b & Quatro extensômetros, sem R0: R2 e R3 em tração, R1 e R4 em compressão (ou o inverso). Resulta V\_o = V\_s G $\varepsilon$, e $\alpha$$\Delta$T se cancela & Quatro extensômetros: R2 e R3 na tração, R1 e R4 na compressão (ou tudo trocado), com V0 = Vs G $\varepsilon$, sem efeito térmico & igual & ambos \\
\bottomrule\end{longtable}%
\begin{itemize}%
\item%
\textbf{a)} Pelo divisor de tensão, V\_o $\approx$ (V\_s/4)(x2+x3-x1-x4) com x\_i = $\Delta$R\_i/R0. Com só um extensômetro, |V\_o| = V\_s|x|/4, com x = G$\varepsilon$+$\alpha$$\Delta$T. Conferi numericamente, para G=2, $\varepsilon$=1e{-}4, $\alpha$$\Delta$T=3e{-}5: exato 5,7493e{-}5$\cdot$V\_s contra linear 5,75e{-}5$\cdot$V\_s, nos quatro braços (sinais -,+,+,- para R1..R4). A figura do PDF confere com a transcrição.%
\item%
\textbf{b)} Com x2=x3=G$\varepsilon$+$\alpha$$\Delta$T e x1=x4=-G$\varepsilon$+$\alpha$$\Delta$T, V\_o = (V\_s/4)$\cdot$4G$\varepsilon$ = V\_sG$\varepsilon$. A conta exata dá V\_sG$\varepsilon$/(1+$\alpha$$\Delta$T). Numericamente: 1,99994e{-}4 contra 2e{-}4, e a diferença é um termo $\varepsilon$$\alpha$$\Delta$T, desprezível pelo enunciado. A resolução ainda justifica a otimalidade pela soma dos s\_i, que vale no máximo 4, o que é coerente com o gabarito.%
\end{itemize}%
\paragraph{Pontos de dúvida}%
\begin{itemize}%
\item%
\textit{(statement)} A polaridade da fonte vem do símbolo da bateria (placa maior em cima, então A positivo). Ela só muda o sinal de V\_o, sem alterar o módulo pedido em a) nem a montagem de b). Tanto a resolução quanto o gabarito adotam V\_A = V\_s.%
\item%
\textit{(statement)} Em b), 'todo o sistema está na mesma temperatura' é lido como 'todos os dispositivos sofrem o mesmo $\Delta$T', e não como $\Delta$T = 0. As duas resoluções usam essa leitura, e com ela o cancelamento térmico só ocorre com extensômetros em braços adjacentes.%
\end{itemize}%
\section*{Questão 8}%
\href{https://poliedroresolve.sistemapoliedro.com.br/vestibulares/ita/2027/ita-2-fase-dia-3-22-10-2026-2-fase-dia-3/dissertativa-8-fisica-2-fase-dia-3-ita-2027}{Resolução do gabarito}\par\smallskip%
\begin{longtable}{@{}p{0.06\textwidth} p{0.33\textwidth} p{0.33\textwidth} p{0.1\textwidth} p{0.08\textwidth}@{}}\toprule Item & Resolução & Gabarito & Veredito & Certo \\ \midrule
a & v = $\surd$(2qV/m) & v = $\surd$(2qV/m) & igual & ambos \\
b & Movimento helicoidal uniforme: MRU ao longo de B com v cos$\theta$ e MCU no plano perpendicular, com R = mv sen$\theta$/(qB) e T = 2$\pi$m/(qB). A justificativa é que F = qv$\times$B é perpendicular a B (não há força ao longo de B) e perpendicular a v (não realiza trabalho, então |v| é constante). & Movimento helicoidal uniforme: MCU no plano perpendicular a B (R = mv sen$\theta$/(qB), T = 2$\pi$m/(qB)) composto com MU ao longo de B (v cos$\theta$) & igual & ambos \\
c & cos$\theta$ = (LB/(2$\pi$n))$\cdot$$\surd$(q/(2mV)), n = 1, 2, 3, \ldots{} (com cos$\theta$ < 1) & cos$\theta$ = qBL/(2$\pi$nmv) = (BL/(2$\pi$n))$\cdot$$\surd$(q/(2mV)), n $\in$ $\mathbb{N}$* (com cos$\theta$ $\leq$ 1) & igual & ambos \\
\bottomrule\end{longtable}%
\begin{itemize}%
\item%
\textbf{a)} Teorema trabalho--energia cinética, qV = mv$^2$/2, saindo do repouso. A resolução e o gabarito dão o mesmo resultado.%
\item%
\textbf{b)} Conclusão, justificativa, raio e período iguais aos do gabarito. A resolução também diz que v? $\neq$ 0 e v? $\neq$ 0 porque 0 < $\theta$ < 90$^\circ$ na figura, o que garante que a trajetória é mesmo uma hélice.%
\item%
\textbf{c)} Refiz a conta. A partícula só volta à linha de campo de 1 depois de n períodos completos, e então L = n$\cdot$v cos$\theta$$\cdot$2$\pi$m/(qB). Disso sai cos$\theta$ = qBL/(2$\pi$nmv) e, com v do item a, cos$\theta$ = (LB/2$\pi$n)$\surd$(q/2mV). Conferi numericamente em Python, com valores aleatórios, que as duas formas são iguais e que n$\cdot$v cos$\theta$$\cdot$T = L. A única diferença é a desigualdade estrita na restrição (cos$\theta$ < 1 contra cos$\theta$ $\leq$ 1). Ela só afeta o caso degenerado $\theta$ = 0, em que não há hélice, e não altera a resposta.%
\end{itemize}%
\paragraph{Pontos de dúvida}%
\begin{itemize}%
\item%
\textit{(solution)} Para os n admissíveis, a resolução usa desigualdade estrita (cos$\theta$ < 1, n > (LB/2$\pi$)$\surd$(q/2mV)) e o gabarito usa n $\geq$ \ldots{} (cos$\theta$ $\leq$ 1). A igualdade corresponde a $\theta$ = 0: a partícula anda sobre a própria linha de campo, sem hélice e sem 'retorno' propriamente dito. Excluir esse caso, como faz a resolução, é coerente com a figura ($\theta$ > 0), e nenhuma das duas formas deve mudar a nota.%
\item%
\textit{(solution)} Erro de redação: 'a partícula descendo reta pela linha de campo' no caso limite $\theta$ $\to$ 0. O certo seria 'movendo{-}se em linha reta ao longo da linha de campo'. Não afeta a física.%
\item%
\textit{(statement)} O item c pede 'a condição que cos$\theta$ deve satisfazer'. A resposta esperada é uma família discreta, cos$\theta$ = (LB/2$\pi$n)$\surd$(q/2mV) com n inteiro positivo, e não um único valor. A resolução e o gabarito interpretam da mesma forma, deixando n livre e lembrando a restrição cos$\theta$ $\leq$ 1.%
\end{itemize}%
\section*{Questão 9}%
\href{https://poliedroresolve.sistemapoliedro.com.br/vestibulares/ita/2027/ita-2-fase-dia-3-22-10-2026-2-fase-dia-3/dissertativa-9-fisica-2-fase-dia-3-ita-2027}{Resolução do gabarito}\par\smallskip%
\begin{longtable}{@{}p{0.06\textwidth} p{0.33\textwidth} p{0.33\textwidth} p{0.1\textwidth} p{0.08\textwidth}@{}}\toprule Item & Resolução & Gabarito & Veredito & Certo \\ \midrule
a & E = E0{[}1+(-1)\^{}n{]} cos(kx - $\omega$t) $\hat z$: onda progressiva no sentido +x (velocidade de fase $\omega$/k = c > 0), com amplitude E0{[}1+(-1)\^{}n{]}, que vale 2E0 para n par e 0 para n ímpar. & E = E0{[}1 + (-1)\^{}n{]} cos(kx - $\omega$t) $\hat z$; amplitude E0{[}1 + (-1)\^{}n{]}, onda no sentido +x & igual & ambos \\
b & n par (d = m$\lambda$): interferência construtiva total, com amplitude 2E0. n ímpar (d = $\lambda$/2, 3$\lambda$/2, \ldots{}): interferência destrutiva total, com campo nulo. & n par: construtiva (amplitude 2E0); n ímpar: destrutiva total (amplitude 0) & igual & ambos \\
c & B = -B0 cos(kx - $\omega$t) $\hat y$, com B0 = 2E0/c para n par (B aponta em -$\hat y$ quando E aponta em +$\hat z$, e E$\times$B fica em +x?); B = 0 para n ímpar. & B = -B0 cos(kx - $\omega$t) $\hat y$ (E $\times$ B no sentido +x?) & igual & ambos \\
\bottomrule\end{longtable}%
\begin{itemize}%
\item%
\textbf{a)} Como k(x+d) = kx + n$\pi$ e cos($\theta$+n$\pi$) = (-1)\^{}n cos$\theta$, as duas chegam à mesma expressão. A forma alternativa 2E0 cos(n$\pi$/2) cos(kx - $\omega$t + n$\pi$/2), que a resolução cita, coincide com essa para todo n inteiro (para n = 2: (-1)$\cdot$(-cos$\theta$) = cos$\theta$). O argumento sobre o sentido de propagação está correto: a fase constante dá x = ($\omega$/k)t + const.%
\item%
\textbf{b)} As conclusões são as mesmas. A resolução ainda traduz a condição como d = n$\lambda$/2 e observa que a intensidade fica 4 vezes maior, o que está correto e é um complemento.%
\item%
\textbf{c)} A expressão e a orientação são as mesmas. A resolução deduz o campo pela lei de Faraday: $\partial$B\_y/$\partial$t = $\partial$E\_z/$\partial$x = -kE\_m sen(kx-$\omega$t), o que dá B\_y = -(k/$\omega$)E\_m cos(kx-$\omega$t) = -(E\_m/c)cos(kx-$\omega$t). Refiz a conta e confere. O gabarito chega ao mesmo resultado só pela condição $\hat z$ $\times$ (-$\hat y$) = x?. A resolução também identifica B0 = 2E0/c (para n par), o que é um complemento coerente.%
\end{itemize}%
\paragraph{Pontos de dúvida}%
\begin{itemize}%
\item%
\textit{(statement)} O enunciado diz que d é uma constante positiva, mas inclui n = 0 (d = 0) na condição kd = n$\pi$. Isso não muda nenhum item, mas a resolução usa n = 0 na conferência.%
\item%
\textit{(statement)} No item c, o enunciado só diz que B tem 'amplitude de módulo B0'. Para n ímpar o campo resultante é nulo, então essa amplitude só faz sentido (B0 $\neq$ 0) para n par. A resolução diz isso explicitamente e acrescenta B0 = 2E0/c, que o gabarito não pede. A resposta não muda, só fica mais completa.%
\item%
\textit{(solution)} A resolução também discute a leitura alternativa em que B0 seria a amplitude de cada onda componente (B = -B0{[}1+(-1)\^{}n{]}cos(kx-$\omega$t)$\hat y$). Isso não contradiz a resposta principal, que coincide com a do gabarito, mas um corretor pode achar que a resposta ficou com mais de uma versão. A resposta destacada no quadro é a do gabarito.%
\end{itemize}%
\section*{Questão 10}%
\href{https://poliedroresolve.sistemapoliedro.com.br/vestibulares/ita/2027/ita-2-fase-dia-3-22-10-2026-2-fase-dia-3/dissertativa-10-fisica-2-fase-dia-3-ita-2027}{Resolução do gabarito}\par\smallskip%
\begin{longtable}{@{}p{0.06\textwidth} p{0.33\textwidth} p{0.33\textwidth} p{0.1\textwidth} p{0.08\textwidth}@{}}\toprule Item & Resolução & Gabarito & Veredito & Certo \\ \midrule
a & $\nu$' = $\nu$ / {[}1 + (h$\nu$/(m\_e c$^2$))(1 - cos $\theta$){]} (efeito Compton), deduzido pela conservação de energia e momento relativísticos & $\nu$' = $\nu$/(1 + (h$\nu$/(m\_e c$^2$))(1 - cos$\theta$)) (Compton) & igual & ambos \\
b & V = -c($\nu$$^2$ - $\nu$'$^2$)/($\nu$$^2$ + $\nu$'$^2$), com V > 0 no sentido do fóton espalhado. Ou seja, |V| = c$\cdot$$\varepsilon$(2+$\varepsilon$)/(2+2$\varepsilon$+$\varepsilon$$^2$), com $\varepsilon$ = (h$\nu$/(m\_e c$^2$))(1 - cos $\theta$), e O se move sobre a reta do fóton espalhado em sentido oposto ao dele (indo ao encontro do fóton) & O se move na direção do fóton espalhado, em sentido oposto ao da propagação. $\beta$ = ($\nu$$^2$ - $\nu$'$^2$)/($\nu$$^2$ + $\nu$'$^2$), ou seja, V = c$\cdot${[}(1+$\varepsilon$)$^2$ - 1{]}/{[}(1+$\varepsilon$)$^2$ + 1{]}, com $\varepsilon$ = h$\nu$(1 - cos$\theta$)/(m\_e c$^2$) & equivalente & ambos \\
\bottomrule\end{longtable}%
\begin{itemize}%
\item%
\textbf{a)} Expressão idêntica e dedução igual: conservação do momento (lei dos cossenos), conservação da energia e E$^2$ = p$^2$c$^2$ + m$^2$c$^{4}$, chegando a m\_e c$^2$($\nu$ - $\nu$') = h$\nu$$\nu$'(1 - cos $\theta$). Os casos{-}limite conferem ($\theta$ = 0 dá $\nu$' = $\nu$; $\nu$' < $\nu$ quando $\theta$ $\neq$ 0).%
\item%
\textbf{b)} Refiz a conta. Do Doppler longitudinal, $\nu$'' = $\nu$'$\surd$((1+$\beta$)/(1-$\beta$)) quando O se aproxima do fóton. Impondo $\nu$'' = $\nu$, sai $\beta$ = ($\nu$$^2$ - $\nu$'$^2$)/($\nu$$^2$ + $\nu$'$^2$). Com $\nu$/$\nu$' = 1+$\varepsilon$, isso dá $\beta$ = ((1+$\varepsilon$)$^2$ - 1)/((1+$\varepsilon$)$^2$ + 1) = $\varepsilon$(2+$\varepsilon$)/(2+2$\varepsilon$+$\varepsilon$$^2$), mesma expressão nas duas formas. Conferi numericamente para $\varepsilon$ = 0,01; 0,1957; 1,3 e 4: as três expressões coincidem, e o boost devolve $\nu$'' = $\nu$. O módulo e o sentido físico (aproximar{-}se do fóton, para ter desvio para o azul) são os mesmos. Só muda a convenção: a resolução escreve V com sinal negativo (eixo no sentido do fóton), o gabarito dá o módulo e diz o sentido por extenso. A aproximação extra para ângulo pequeno da resolução, |V| $\approx$ h$\nu$$\theta$$^2$/(2m\_e c), é correta (vale para $\theta$ ? 1 e $\varepsilon$ ? 1) e não contradiz nada.%
\end{itemize}%
\paragraph{Pontos de dúvida}%
\begin{itemize}%
\item%
\textit{(solution)} Convenção de sinal no item b. A resolução apresenta V = -c($\nu$$^2$ - $\nu$'$^2$)/($\nu$$^2$ + $\nu$'$^2$) (negativo porque o eixo aponta no sentido do fóton espalhado). O gabarito dá V positivo e indica o sentido por extenso ("sentido oposto ao de propagação"). O conteúdo físico é o mesmo, mas convém o professor confirmar que o sinal negativo, explicado na própria resolução, não seria penalizado.%
\item%
\textit{(key)} A resolução do gabarito tem erros de digitação ou de transcrição que não mudam a resposta. Escreve p\_e$^2$ = h$\nu$/c$^2$ + (h$\nu$')/c$^2$ - ..., quando o certo é (h$\nu$/c)$^2$ + (h$\nu$'/c)$^2$ - ...; e escreve (1 + h$\nu$/(m\_ec$^2$)(1-cos$\theta$)\^{}2 - 1), em que o quadrado deveria cobrir todo o (1 + $\varepsilon$) e não só (1 - cos$\theta$). A resposta por item do gabarito, V = c{[}(1+$\varepsilon$)$^2$ - 1{]}/{[}(1+$\varepsilon$)$^2$ + 1{]}, está correta.%
\item%
\textit{(statement)} "Mesma direção do fóton espalhado" é ambíguo entre direção (a reta) e sentido. Se fosse no mesmo sentido do fóton, não haveria solução, porque O veria uma frequência ainda menor que $\nu$'. Tanto o gabarito quanto a resolução entendem "direção" como a reta e concluem que O deve ir em sentido oposto ao do fóton.%
\end{itemize}%
\end{document}